\section{Tokenização e arquitetura}
\label{sec:arquitetura}
\subsection{Duas unidades de entrada}
\textbf{Palavras.} Convertemos o prefixo limpo para minúsculas. A expressão
\verb|[a-z]+(?:'[a-z]+)?|\verb!|\d+! preserva palavras, contrações e números.
O vocabulário contém os 19.998 itens mais frequentes no treino, além de
\texttt{<pad>} (ID 0) e \texttt{<unk>} (ID 1). A frequência determina a ordem,
com desempate lexical. A posição de cada ID é preservada na sequência.
Palavras de negação como \textit{not} seguem para o modelo.

\textbf{BPE.} Mantemos a normalização lexical e separamos palavras com
\texttt{WhitespaceSplit}. O BPE de caracteres aprende uniões de pares
frequentes com os 6.144 textos de treino, frequência mínima 2 e vocabulário
de 20.000 itens. O marcador \texttt{</w>} identifica o fim de palavra.
O tokenizer é congelado antes do treino supervisionado da CNN.
O procedimento adapta BPE para unidades de texto \cite{sennrich}.

\caso{Entrada}{\ExemploCobertura}
\begin{tabularx}{\linewidth}{@{}lX@{}}
\toprule
\tabelahead{Palavras} & \texttt{the \enspace <unk> \enspace scene \enspace was \enspace unexpectedly \enspace moving}\\
\tabelahead{BPE} & \texttt{the</w> \enspace p \enspace all \enspace bearer</w> \enspace scene</w>} \newline
\texttt{was</w> \enspace unexpectedly</w> \enspace moving</w>}\\
\bottomrule
\end{tabularx}
\legendinha{Tabela 3. Saída dos tokenizadores congelados. \texttt{pallbearer} tem
um ID desconhecido com palavras e três peças com BPE.}

\subsection{Convolução sobre embeddings}
Cada ID acessa um vetor de 128 dimensões aprendido com a tarefa. O filtro
de largura \(k\) lê \(k\) tokens vizinhos e usa os mesmos pesos em todas as
posições. Aplicamos três grupos paralelos de convoluções de larguras 3, 4 e
5, com cem filtros por largura. A ReLU deixa respostas positivas passarem.
O máximo global seleciona a maior ativação por filtro. A topologia segue
a família de CNNs para classificação de sentenças \cite{kim,prince}.

\vspace{2mm}
\Arquitetura
\legendinha{Figura 1. TextCNN do projeto. \(B\) é o tamanho do minibatch e \(T\) é
o comprimento após padding. A camada linear produz logits negativo e positivo.}

Se \(e_t\) é o embedding na posição \(t\), o filtro \(f\) resume a entrada por
\[
  m_f = \max_{t\in\mathcal V_k}
  \operatorname{ReLU}\!\left(\langle W_f, [e_t,\ldots,e_{t+k-1}]\rangle + b_f\right).
\]
\(\mathcal V_k\) exclui as janelas que avançam sobre padding em entradas com
pelo menos \(k\) tokens. Para sequências mais curtas, a implementação mantém
a primeira janela com padding. A concatenação dos 300 máximos passa por dropout de
0,5 no treino e uma camada \(\mathrm{Linear}(300,2)\).

Na CNN principal, \(20.000\times128=2.560.000\) parâmetros pertencem aos
embeddings. As convoluções somam \(100\sum_{k\in\{3,4,5\}}(128k+1)=153.900\).
A saída soma \(2(300+1)=602\), totalizando \textbf{2.714.502 parâmetros}.
A mesma largura em tokens pode abranger menos palavras quando BPE fragmenta
uma palavra. \fonte{src/imdb_lab/models.py}{Implementação}.
